\emph{Översikt}
Listan är en behållare som kan ändras.  Den här modulen inför dess
oföränderliga syskon, \emph{tupeln}: en följd av värden som hör ihop och
som inte ska kunna ändras efteråt --- en koordinat, ett fullständigt
namn, ett par av ett nummer och det som numret hör till.  Vi ser vad som
skiljer en tupel från en lista genom att göra samma sak med båda, och vi
använder \emph{uppackning} för att plocka isär en tupel i flera
variabler.  Det gör det möjligt för en funktion att lämna tillbaka flera
värden på en gång, och det är den formen
\mintinline{python}{enumerate} har när vi vill ha både numret och
elementet ur en behållare.

\emph{Lärandemål}
Efter denna modul ska studenten kunna:
\begin{restatable}{lo}{TupleLOCreate}\label{TupleLOCreate}%
  Skapa en tupel, plocka ut ett element ur den med index och gå igenom
  den med \mintinline{python}{for}.
\end{restatable}% Lo09
\begin{restatable}{lo}{TupleLOImmutable}\label{TupleLOImmutable}%
  Förklara vad som skiljer en tupel från en lista, och välja tupel när
  värdena hör ihop och inte ska ändras var för sig.
\end{restatable}% Lo09
\begin{restatable}{lo}{TupleLOUnpack}\label{TupleLOUnpack}%
  Packa upp en tupel i flera variabler, och använda det för att låta en
  funktion lämna tillbaka flera värden.
\end{restatable}% Lo09, Lo03
\begin{restatable}{lo}{TupleLOEnumerate}\label{TupleLOEnumerate}%
  Gå igenom en behållare och få både elementets nummer och elementet,
  med \mintinline{python}{enumerate}, och förklara vad
  \mintinline{python}{enumerate} lämnar tillbaka.
\end{restatable}% Lo04
\begin{restatable}{lo}{TupleLODocs}\label{TupleLODocs}%
  Hitta och läsa dokumentationen för tupler och för
  \mintinline{python}{enumerate}, och därifrån svara på en fråga som
  inte besvarats på föreläsningen.
\end{restatable}% Lo09

\emph{Förkunskaper}
Modulen förutsätter listor (index, \mintinline{python}{len} och
\mintinline{python}{for} över en lista) samt egna funktioner med
parametrar och returvärde.

\emph{Läsning}
Kursens videogenomgång \emph{Tupler} täcker samma material.
